\documentclass[11pt]{article}
\usepackage[a4paper,margin=27mm]{geometry}
\usepackage[T1]{fontenc}
\usepackage{lmodern,amsmath,amssymb,amsthm,microtype,booktabs}
\usepackage[colorlinks=true,allcolors=blue]{hyperref}
\newtheorem{theorem}{Theorem}
\newtheorem{lemma}[theorem]{Lemma}
\newtheorem{proposition}[theorem]{Proposition}
\newcommand{\Tr}{\operatorname{Tr}}
\newcommand{\id}{\operatorname{id}}
\newcommand{\CEA}{C_{\mathrm{EA}}}
\newcommand{\cmp}{\mathrm{cmp}}
\newcommand{\PhiT}{\Phi_{\kappa,n}}
\title{Sharp entropy-budget laws for Gaussian reset channels}
\author{Lluis Eriksson}
\date{26 October 2026\\\small Working draft 663}
\begin{document}
\maketitle
\begin{abstract}
An unflagged reset channel transmits through a Gaussian signal channel
with probability \(\lambda\), and otherwise emits a fixed state whose
entropy, but not energy or orientation, is bounded. We determine the
complete high-energy minimax loss curves for unassisted classical
communication, permitting arbitrary entangled block inputs and collective
decoding. Relative to \(\lambda\) times the signal's own classical
capacity, the absolute loss is
\(b(\lambda,(1-\lambda)t_\alpha)+o(1)\), uniformly in every finite entropy
budget. For Shannon entropy
\(t_1=\min\{1,B/\log_2(E+1)\}\); for fixed finite \(\alpha>1\),
\(t_\alpha=1-2^{-(\alpha-1)B/\alpha}\).
The sharp regime is signal vacuum noise \(t=o(E)\).
It covers thermal attenuation and amplification, and additive Gaussian
noise bounded away from zero. A larger unknown fixed family has the
same law and a common encoder/decoder, including cold and noiseless
members. No preshared entanglement or additivity of the reset mixture
is assumed. The converse uses causal quantum references, a summed
Gaussian minimum-output-entropy floor and averaged Choi-order deficits.
Selecting a fraction of good positions and using established telescopic
relative entropy gives an error uniform in block length.
Finite cq compound coding with exact physical types supplies the
matching operational rate for arbitrary reset energy. A strict flag-loss
gap at bounded signal information proves necessity.
The assisted thermal law is also proved, relative to its own benchmark.
The actual critical loss curve and efficient blocklength bounds remain
open here. Working-argument audits support the thermal ingredients;
the enlarged exact manuscript requires its own full assessment.
\end{abstract}

\section{Model, scope and theorem}
All logarithms in information quantities are base two; \(\ln\) is natural.
Let \(H=\sum_{j\ge0}j|j\rangle\langle j|\),
\(g(x)=(x+1)\log_2(x+1)-x\log_2x\), and let \(h\) denote binary entropy.
The thermal state of mean \(x\) is \(\gamma_x\).
A thermal attenuator \(\PhiT\) mixes the input with an environment thermal
state of mean \(n\), through a beamsplitter of photon transmissivity
\(0<\kappa<1\). Write \(\delta=1-\kappa\), \(q=1-\lambda\), with
\(0<\lambda<1\), and
\[
 {\cal N}_{\kappa,n,\tau}(\rho)=\lambda\PhiT(\rho)+q\tau\Tr\rho.
\]
The reset \(\tau\) has no energy or orientation restriction.
All parameters are fixed throughout a code block. The input constraint
is expected physical block-average energy at most \(E\), not peak energy.
For the assisted baseline developed first, entanglement assistance is
unlimited. Sections~\ref{sec:unassisted}--\ref{sec:generalU} remove that
resource and permit arbitrary entangled block encoding without preshared
entanglement. Classical error is average-message
error. In a compound capacity one encoder and decoder serve the whole
fixed-parameter family, with uniform error. The coding limit is taken
at each fixed \(E,n,B\), before the energy limit.

Put \({\cal T}_{\alpha,B}=\{\tau:S_\alpha(\tau)\le B\}\), where
\(S_1=S\) and \(S_\alpha=(1-\alpha)^{-1}\log_2\Tr\tau^\alpha\).
The established constrained Gaussian assisted capacity
\cite{hw,hs} is
\begin{align}
 C_{\kappa,n}(E)&=g(E)+g(M)-g(u_l)-g(u_s),&
 M&=\kappa E+\delta n,\tag{C}\label{eq:C}\\
 V&=\delta E+(1+\kappa)n+1,&
 D&=\sqrt{V^2-4\kappa n(n+1)},\nonumber\\
 u_l&=\tfrac12[D+\delta(E-n)-1],&
 u_s&=\tfrac12[D-\delta(E-n)-1].\nonumber
\end{align}
Here \(g(u_l)+g(u_s)\) is the entropy of the \emph{true two-mode}
complement, including purification of the initial thermal environment.
A one-mode weak complement cannot replace it.

Define the known-reset and unknown-reset deficits
\[
\begin{split}
 D_{\rm ind}(E,n,B)&=\lambda C_{\kappa,n}(E)-
       \inf_{\tau\in{\cal T}_{\alpha,B}}\CEA({\cal N}_{\kappa,n,\tau},H,E),\\
 D_{\cmp}(E,n,B)&=\lambda C_{\kappa,n}(E)-
       \CEA^{\cmp}(\{{\cal N}_{\kappa,n,\tau}:\tau\in{\cal T}_{\alpha,B}\},H,E).
\end{split}
\]
Thus \(D_{\rm ind}\le D_{\cmp}\). Let
\[
 b(u,v)=(u+v)h\!\left(\frac{u}{u+v}\right),\quad f(a)=b(\lambda,qa),
\]
\[
 t_\alpha(E,B)=
 \begin{cases}
 \min\{1,B/\log_2(E+1)\},&\alpha=1,\\
 1-2^{-((\alpha-1)/\alpha)B},&1<\alpha<\infty .
 \end{cases}
\]
Use continuous boundary values in \(b\).
\begin{theorem}[Growing-noise law]\label{thm:main}
Fix \(\kappa,\lambda\in(0,1)\) and \(\alpha=1\) or a fixed finite
\(\alpha>1\). For every prescribed path \(n(E)>0\) with \(n(E)/E\to0\),
\[
 \sup_{0\le B<\infty}\max_{Y\in\{{\rm ind},\cmp\}}
 \left|D_Y(E,n(E),B)-f(t_\alpha(E,B))\right|\longrightarrow0.
 \tag{P}\label{eq:P}
\]
The argument below uses the scientific dependencies disclosed at the end.
It claims neither a uniform blocklength nor uniformity over
all asymptotic paths at a fixed energy without a common envelope.
\end{theorem}
This removes a sufficient \(n=o(\sqrt E)\) restriction of an earlier
working proof, rather than establishing that restriction as a physical
threshold. The proposed addition concerns operational quantifiers and
uniform errors. Gaussian extremality, entropy mixing, Helstrom testing,
telescopic relative entropy and compound coding are prior ingredients.

\begin{theorem}[Sharp regime for the complete curves]\label{thm:sharp}
Fix \(\kappa,\lambda\in(0,1)\), and either \(\alpha=1\) or a fixed finite
\(\alpha>1\). For every finite nonnegative path \(n(E)\), the complete
uniform-in-budget law \eqref{eq:P} holds if and only if \(n(E)=o(E)\).
For the unknown fixed signal family
\(\kappa\in[\kappa_0,1]\), \(n\in[0,N(E)]\), with fixed
\(0<\kappa_0<1\) and finite \(N(E)\ge0\), its compound analogue relative
to \(\lambda C_{\kappa_0,N(E)}(E)\) holds if and only if \(N(E)=o(E)\).
There is no smoothness or monotonicity assumption on either path.
\end{theorem}
Sufficiency is Theorem~\ref{thm:main} and
Proposition~\ref{prop:unknown}; zero-temperature subsequences use the
previous pure-loss law \cite{previous}.
Necessity, proved below, excludes the complete curve on a growing
budget. It does not assert necessity for every restricted small-budget
problem or give the actual critical-temperature loss curve.

\section{A uniform Gaussian entropy estimate}
The input/output energy bounds are
\(S(\rho)\le g(E)\), \(S(\PhiT(\rho))\le g(M)\).
For input mean \(x\), the true complementary entropy is bounded by
\(g(\delta x+\kappa n)+g(n)\).
Let \(\beta_x=\log_2(1+1/x)\) for \(x>0\).
The complementary Gibbs logarithm pulled back to the input is a
quadratic form \(a(E;n)H+b(E;n)I\), with
\[
 a(E;n)=\delta\left[\beta_{u_l}
           +(\beta_{u_l}+\beta_{u_s})r\right],
 \qquad r=\frac12(V/D-1).
 \label{eq:a}
\]
For clarity, a direct verification uses the environment modes
\(e'=-\sqrt\delta\,a+\sqrt\kappa\,e\) and purifier \(z\).
With \(\cosh(2s)=V/D\), \(\sinh(2s)=2\sqrt{\kappa n(n+1)}/D\),
the thermal normal modes are
\(x=\cosh(s)e'-\sinh(s)z^\dagger\) and
\(y=\cosh(s)z-\sinh(s)e'^\dagger\).
The logarithm is a constant plus
\(\beta_{u_l}x^\dagger x+\beta_{u_s}y^\dagger y\).
Pullback of \(e'^\dagger e'\), \(z^\dagger z\), and
\(e'z+e'^\dagger z^\dagger\) is respectively
\(\delta H+\kappa nI\), \(nI\), and
\(2\sqrt{\kappa n(n+1)}I\), giving the displayed coefficient.
These are positive-form/finite-mean expectation identities on the number
form domain. They are not bounded-adjoint identities for an unbounded
logarithm. Finite Fock cutoffs and finite energy justify their limiting
expectations. Differentiating the Gibbs cross entropy at fixed \(n\)
also gives \(a(E;n)=\partial_E[g(u_l)+g(u_s)]\).

Relative-entropy data processing with reference \(\gamma_E\) gives
\[
 S(\rho)-S(\PhiT^c(\rho))
 \le g(E)-g(u_l)-g(u_s)+\zeta,\qquad
 \zeta=E|\beta_E-a(E;n)|.
 \tag{1}
\]
No optimality of thermal reverse coherent information is assumed.
For \(0<n\le E\), the exact radical obeys
\[
 D\ge\delta(E+n)+1,\quad
 \delta E\le u_l\le\delta E+\kappa n,\quad
 \delta n\le u_s\le n,\quad
 0\le r\le\frac{\kappa n(n+1)}{\delta^2E^2}.
\]
The elementary inequalities
\(0\le x^{-1}-\ln(1+x^{-1})\le(2x^2)^{-1}\) imply
\[
 \zeta\le\Xi(E,n):=
 \frac{1+1/\delta}{2E\ln2}
 +\frac{\kappa n}{\delta E\ln2}
 +\frac{2\kappa(n+1)}{\delta^2E\ln2}.
 \tag{2}
\]
Thus \(\Xi=o(1)\) throughout \(n=o(E)\).
All derivatives hold \(n\) fixed; \(n'(E)\) is never used.

\section{Flat Choi order and near-capacity inputs}
\begin{lemma}[Known bounded Choi form, applied here]\label{lem:flat}
For the reference/output state
\(\omega=(\id\otimes\PhiT)(\psi_\rho)\) of the canonical purification,
\[
 \omega\le c\,\rho_R\otimes I_B,\qquad
 c=\frac1{\delta(n+1)}.                                      \tag{3}
\]
\end{lemma}
This bound is an exact specialization of Holevo's established Gaussian
Choi norm formula \cite[Proposition~2 and the one-mode example, p.~13]{choi}.
His amplitude parameter is \(\sqrt\kappa\) and his noise coefficient is
\(m=\delta(n+1/2)\), so the norm is
\([m+|\kappa-1|/2]^{-1}=c\). The following regularization argument
also explains the application to the present purification; the Choi
norm formula is not a new lemma of this paper.
\begin{proof}
Use an auxiliary thermal input of mean \(x>0\), keeping \(n\) fixed.
Its reference/output state \(\omega_x\) has the same nonzero spectrum as
the true Gaussian complement. Formula \eqref{eq:C} gives
\[
 (u_l+1)(u_s+1)
 =\delta xn+\frac{\delta(x+n)}2+\frac{D+1}2
 \ge\delta(n+1)(x+1).
\]
If \(J\) is the canonical positive Choi form and
\(G_x=(x/(x+1))^{H/2}\), then
\((G_x\otimes I)J(G_x\otimes I)=(x+1)\omega_x\le cI\).
For a vector with finite input-number support, insert \(G_x^{-1}\);
its norm converges to the original norm as \(x\to\infty\).
The positive form therefore has bound \(cI\) on a dense domain and
extends to a bounded positive operator. Sandwiching with
\(\sqrt{\rho^T}\otimes I\) gives (3), first for finite truncations and
then in trace norm. No unnormalizable Bell vector is used as a state.
\end{proof}

Put \(C=C_{\kappa,n}(E)\). At fixed \(n\),
\[
 0\le\partial_E C\le\frac2{E\ln2}.                            \tag{4}
\]
The upper bound follows from
\(\partial_E C=\beta_E+\kappa\beta_M-a(E;n)\) and \(a\ge0\);
monotonicity follows from enlargement of the feasible input set.

If \(w=\|\rho\|_\infty\), isolate a largest pure eigencomponent.
Mutual information of that pure input is zero. The entropy-of-mixture
inequality, whose general \(2h(w)\) version is prior
\cite[Eq.~(11), Section~5.2]{shirokov}, yields
\[
 I(\rho,\PhiT)\le2h(w)+(1-w)C_{\kappa,n}(E/(1-w))
 \le(1-w)C+2h(w)+2(1-w)\log_2\frac1{1-w}.
\]
Consequently
\[
 I(\rho,\PhiT)\ge C-\Delta
 \quad\Longrightarrow\quad
 w\le\frac{\Delta+2+2/(e\ln2)}C.                              \tag{5}
\]
The endpoint \(w=1\) is excluded eventually by \(C>\Delta\).
Let \(S=c(M+1)\). Since \(u_l\le E\), \(u_s\le n\) for \(n\le E\),
\[
 C\ge g(M)-g(n)\ge\log_2 S+\log_2\delta-\log_2 e.
 \tag{6}
\]
Thus \(C\to\infty\) whenever \(n=o(E)\).
The established regular-channel identity \(\PhiT(I)=I/\kappa\)
\cite[Proposition~2]{gain} gives
\[
 \Pr_{\PhiT(\rho)}(0)\le w/\kappa=o(1)                       \tag{7}
\]
on the set in (5).

\section{Pinching and branch discrimination}
Let \(\tau_a=(1-a)|0\rangle\langle0|+a\gamma_M\), \(0\le a\le1\).
Flagging gives \(I(\rho,{\cal N}_{\kappa,n,\tau_a})
\le\lambda I(\rho,\PhiT)\).
Inputs below \(\lambda C-h(\lambda)-1\) already satisfy the needed
converse because \(f(a)\le h(\lambda)\). All remaining inputs satisfy
(5) with \(\Delta=(h(\lambda)+1)/\lambda\), independent of \(E,n,a\).

Set \(\sigma_0=\rho_R\otimes I_B\), an \emph{unnormalized}
infinite-trace operator, and
\(D_0(\omega)=\Tr\omega(\log_2\omega-\log_2\sigma_0)
=-H(B|R)_\omega\).
Input/output energy ensures finite entropies and cross entropies.
Equation (3), (6), and \(g(u_l)+g(u_s)\le g(E)+g(n)\) give
\[
 \log_2c-C_0\le D_0(\omega)\le\log_2c,\qquad
 C_0=\Delta+\log_2e-\log_2\delta.                             \tag{8}
\]
We have not applied a normalized-state relative-entropy inequality to
\(\sigma_0\).

For \(S>1\), let \(t=\sqrt{\ln S}\), \(T=\lceil(M+1)t\rceil\),
\(P_B=\sum_{j=1}^T|j\rangle\langle j|\), \(P=I_R\otimes P_B\),
and \(p=\Tr P\omega\). Equations (5),(7) and Markov give
\(p\ge1-w/\kappa-1/t\to1\).
Define normalized good and bad blocks \(\omega_g,\omega_b\).
For these blocks \(D_0\) always uses the original reference \(\rho_R\);
in particular
\[
 D_0(\omega_g)=-H(B|R)_{\omega_g}
                 +D(\rho_R^g\Vert\rho_R).
\]
It is not the negative conditional entropy computed with its own
reference marginal alone.
Pinching preserves the reference marginal and increases entropy by at
most \(h(p)\). Its block formula is
\[
 D_0({\cal P}\omega)=pD_0(\omega_g)+(1-p)D_0(\omega_b)-h(p).
\]
Hence the weighted sum on the right without \(-h(p)\) is at least
\(D_0(\omega)\). By (3),
\(D_0(\omega_b)\le\log_2c-\log_2(1-p)\), and therefore
\[
 D_0(\omega_g)\ge\log_2c-\frac{C_0+h(p)}p.                   \tag{9}
\]
Zero-weight blocks are interpreted by continuity. Reference marginals
are bounded by \(\rho_R/p\), \(\rho_R/(1-p)\), keeping their cross
entropies finite. This step avoids an error of the form \(o(1)\log E\).

Let \(v=\Tr P_B\gamma_M\to1\), \(r_M=M/(M+1)\), and introduce the
\emph{normalized} comparator
\[
 \sigma_g=\rho_R\otimes P_B\gamma_M P_B/v,\qquad
 L=(c/p)v(M+1)r_M^{-T}.
\]
The least thermal eigenvalue on this window gives \(\omega_g\le L\sigma_g\).
Writing \(\mu_g=\Tr H_B\omega_g\ge0\), (9) gives
\[
\begin{split}
 D(\omega_g\Vert\sigma_g)
 &=D_0(\omega_g)+\log_2(M+1)+\beta_M\mu_g+\log_2v\\
 &\ge\log_2S+\log_2v-(C_0+1)/p,\\
 \log_2L&=\log_2S+\log_2v-\log_2p+T\beta_M
          =\log_2S+O(\sqrt{\ln S}).
\end{split}                                                  \tag{10}
\]
In particular the sign of the normalization term is \(+\log_2v\).

For normalized states \(r\le Ls\), Audenaert's established telescopic
inequality \cite[Theorem~4]{audenaert} implies
\[
 D(r\Vert s)\le\frac{\log_2L}{1-1/L}\,
                    \frac12\|r-s\|_1.                      \tag{11}
\]
Indeed write \(s=L^{-1}r+(1-L^{-1})z\) and use that theorem, together
with \(\tfrac12\|r-s\|_1=(1-L^{-1})\tfrac12\|r-z\|_1\).
For infinite reference, common finite compressions with an orthogonal
failure flag preserve normalization and operator order. Their embedded
states converge in trace norm. Lower semicontinuity of relative entropy
and data processing give convergence of the relative entropies, proving
(11) in this case. Equation (10) now forces
\(\frac12\|\omega_g-\sigma_g\|_1\to1\).

The selected proof of Audenaert's Theorem~4 in v2 has a printed sign
inconsistency on p.~10: after defining a difference proportional to
\(\rho-\sigma\), it later uses the opposite orientation.
Taking the auxiliary difference as \((\sigma-\rho)/
(\tfrac12\|\rho-\sigma\|_1)\) restores that displayed identity and
leaves the positive-part norm bound and theorem unchanged.
The use of (11) is the established theorem, not a new inequality.

A Helstrom projection \(Q\le P\) has both normalized testing errors
tending to zero. Its acceptance on the original signal is
\(p\Tr Q\omega_g\to1\). On the original reset it is
\[
 \Tr[Q(\rho_R\otimes\tau_a)]=av\Tr Q\sigma_g\to0,
\]
because the vacuum component is excluded. These bounds are uniform in
\(a\) and in all relevant inputs. Binary Fano gives \(H(F|RB)=o(1)\).
This analytical test may depend on the input; it is a converse device,
not the common operational decoder.

\section{All-input converse and the matching common code}
The exact binary-branch chain rule is
\[
 I(\rho,{\cal N}_{\kappa,n,\tau})=F_\tau(\rho)+H(F|RB),
\quad
 F_\tau=\lambda[S(\rho)-S(\PhiT^c(\rho))]
 +S(\lambda\PhiT(\rho)+q\tau)-qS(\tau)-h(\lambda).
\]
For the finite-entropy witness \(\tau_a\), entropy mixing and (1) give
\[
 F_{\tau_a}\le\lambda C-f(a)+\lambda\Xi
                        +qh\!\left(\frac a{M+1}\right).
 \tag{12}
\]
To see the mixing step, group the signal and thermal reset, of total
weight \(P_a=\lambda+qa\), against the vacuum. The normalized group
has mean at most \(M\), so its entropy is at most \(g(M)\).
Moreover \(S(\tau_a)\ge ag(M)+h(a)-h(a/(M+1))\).
The identity \(h(P_a)-qh(a)-h(\lambda)=-b(\lambda,qa)\)
then gives (12).
Combining (12) with the preceding branch test bounds \emph{every}
admissible input. Evaluation at \(\gamma_E\) supplies the other direction:
\[
 f(a)-o(1)\le
 \lambda C-\CEA({\cal N}_{\kappa,n,\tau_a},H,E)
 \le f(a)+b(\lambda/(M+1),q).                               \tag{13}
\]
Uniformity in \(a\) is essential for changing budgets.

Here is an explicit common error envelope for the converse.
Let \(C_\Delta=\Delta+2+2/(e\ln2)\),
\(H_S=\log_2S+\log_2\delta-\log_2e\), and, eventually,
\[
\begin{aligned}
 w_*&=C_\Delta/H_S,& p_*&=1-w_*/\kappa-1/t,\\
 v_*&=1-(1+\kappa E)^{-1}-e^{-t},&
 K_*&=-\log_2p_*+(2t+1)/\ln2+(C_0+1)/p_*,\\
 \varepsilon_*&=(Sv_*)^{-1}+K_*/\log_2(Sv_*),&
 z_*&=\lambda(1-p_*)+\varepsilon_* .
\end{aligned}
\]
These formulas are used only when \(H_S>0\), \(p_*,v_*>0\),
\(Sv_*>1\), \(M\ge1\), and \(z_*\le1/2\), all eventual
conditions in the joint limit. From (10),(11),
\(1-\tfrac12\|\omega_g-\sigma_g\|_1\le\varepsilon_*\).
Thus a valid error in the left side of (13) is
\[
 R_-=\lambda\Xi(E,n)+qh((1+\kappa E)^{-1})+h(z_*)\longrightarrow0.
 \tag{13a}
\]
This exposes the dependence on \(E,n\), rather than assigning a common
rate to every unspecified sublinear trajectory.

For completeness the achievable bound does not use that witness.
For \(E>1\), let
\[
 L_{\rm in}=\lceil4E\ln(E+1)\rceil,\quad
 \epsilon=(1-1/E)^{L_{\rm in}}\le(E+1)^{-4}.
\]
Truncate \(\gamma_{E-1}\) to its first \(L_{\rm in}\) levels and normalize,
obtaining \(\sigma\) of mean strictly below \(E\).
To expose a uniform bound, write
\(\gamma_{E-1}=(1-\epsilon)\sigma+\epsilon\theta\).
The normalized tail \(\theta\) is a shifted thermal state, with
\(S(\theta)=g(E-1)\). Mutual-information mixing and
\(I(\theta,\PhiT)\le2S(\theta)\) give
\[
 C_{\kappa,n}(E-1)-I(\sigma,\PhiT)
 \le\frac{2h(\epsilon)+2\epsilon g(E-1)}{1-\epsilon}.
\]
Since \(\epsilon<1/16\) for \(E>1\), this is bounded by the larger
entropy error below. Adding (4) at the same \(n\) yields
\[
 I(\sigma,\PhiT)\ge C-U_E,\quad
 U_E=3h(\epsilon)+4\epsilon[g(E)+g(E+L_{\rm in})]
                         +2\log_2\frac E{E-1}=o(1).         \tag{14}
\]
The bound is uniform for \(n\le E\). It compares inputs of energies
\(E-1,E\) at fixed channel, not \(n(E-1)\) with \(n(E)\).

Apply one receiver Fock cutoff \(T_{\rm out}=\lceil E^4\rceil\)
and replace its complement by an orthogonal failure symbol.
The signal failure mass is at most \(E^{-3}\); no energy assumption
is imposed on the reset. With
\[
 e_E=E^{-3/2}+\tfrac12E^{-3},\quad
 W_E=2e_E\log_2L_{\rm in}
       +(1+e_E)h(e_E/(1+e_E)),
\]
finite-reference continuity \cite[Lemma~2]{winter} loses at most
\(W_E=o(1)\). The compressed signal norm is at most
\[
 v_E=\max\left\{\frac1{[1+\kappa(E-1)](1-\epsilon)},E^{-3}\right\}.
\]
A top-\(m\) reset spectral subspace has retained signal mass at most
\((m+1)v_E\), since the failure flag increases rank by at most one.
If \(\eta_m\) bounds its original discarded spectral mass, the binary
observation of that subspace gives
\[
 H(F|B)\le b(\lambda,q\eta_m)
       +b(\min\{\lambda,\lambda(m+1)v_E\},q).
\]
The branch chain rule thus yields
\[
 D_{\cmp}\le f(\eta_m)+R_+(E,m),\quad
 R_+=\lambda(U_E+W_E)+
          b(\min\{\lambda,\lambda(m+1)v_E\},q),               \tag{15}
\]
which vanishes if \(m/E\to0\).

Here is the operational step behind (15).
After the two fixed cuts the channel family has finite input/output
dimension, though its reset index can be arbitrary.
Approximate the diagonal input \(\sigma\) by rational probabilities
while retaining strict mean-cost slack. On the exact type subspace of
a type length \(k_{\rm type}\), every vector has total physical number
\(k_{\rm type}\sum_j jp_j<k_{\rm type}E\).
For completeness, let \(\pi_T\) be the maximally mixed exact-type state
and put \(d_k=D(\pi_T\Vert\sigma_p^{\otimes k})
=kH(p)-S(\pi_T)=O(\log k)\). Its one-site marginals equal \(\sigma_p\).
Data processing bounds the output total correlation by \(d_k\), while
complementary entropy is subadditive. Thus
\[
 I(\pi_T,{\cal M}^{\otimes k})
 \ge kI(\sigma_p,{\cal M})-2d_k
\]
uniformly in every finite channel \({\cal M}\).
This type estimate is also given in \cite[Lemma~3]{previous};
it is distinct from the compound achievability result of
\cite[Lemma~9]{compound}.
Apply the finite-dimensional compound assisted coding theorem
\cite[Theorem~1]{compound} on this subspace. First fix the type and
all cuts for a rate slack, then take the outer coding length to infinity;
only subsequently optimize the finite choices and take \(E\to\infty\).
This preserves physical expected cost and provides a single encoder
and decoder for all resets. It supplies no uniform blocklength or
efficient code. These are uses of established coding results, not
a new coding theorem.

\section{Budget normalization and unknown signal envelopes}
For Shannon entropy,
\(\eta_m\le\min\{1,B/\log_2(m+1)\}\).
Choose \(m=\lfloor E/\ln(E+1)\rfloor\).
Its clipping curve differs uniformly from \(t_1(E,B)\) by \(o(1)\).
For the converse choose
\[
 a_B=\operatorname{clip}_{[0,1]}\frac{B-1}{g(M)}.
\]
The witness is feasible because \(S(\tau_a)\le h(a)+ag(M)\);
the clipped endpoints are feasible separately.
Since \(\kappa E\le M\le E\),
\(\log_2\kappa\le g(M)-\log_2(E+1)\le\log_2e\).
It follows that \(\sup_B|a_B-t_1(E,B)|=o(1)\).
For fixed \(\alpha>1\), choose
\(a_B=1-2^{-((\alpha-1)/\alpha)B}\) and \(m=\lfloor\sqrt E\rfloor\).
The witness obeys \(\Tr\tau_a^\alpha\ge(1-a)^\alpha\), hence is feasible.
The established spectral-tail estimate gives
\(\eta_m\le\min\{1,a_B+(m+1)^{(1-\alpha)/\alpha}\}\).
Since \(f\) is increasing and concave, with
\(|f(x)-f(y)|\le f(|x-y|)\), (13),(15) give (P).
Resets of infinite von Neumann entropy can occur for \(\alpha>1\);
their achievability is evaluated only after finite output compression.
Uncompressed reset entropies are subtracted only for the finite-entropy
witness.

\begin{proposition}[Common code for unknown signal]\label{prop:unknown}
Let known \(0<\kappa_0<1\) and \(N(E)>0\) satisfy \(N(E)=o(E)\).
Suppose \(\kappa\in[\kappa_0,1]\), \(n\in[0,N(E)]\), and
\(\tau\in{\cal T}_{\alpha,B}\) are all unknown and fixed throughout
the block. Then the compound deficit relative to
\(\lambda C_{\kappa_0,N(E)}(E)\) is \(f(t_\alpha(E,B))+o(1)\),
uniformly in finite \(B\). The zero-envelope case is the earlier
pure-loss result, not a substitution in singular thermal logarithms.
\end{proposition}
\begin{proof}
For signal channels only, add receiver Gaussian noise of mean
\((1-\kappa)(N-n)\), raising the environment parameter to \(N\).
Then compose with a thermal attenuator of transmissivity
\(\kappa_0/\kappa\) and environment occupation \(N\).
Its resulting noise is \((1-\kappa_0)N\), so the composition is
\(\Phi_{\kappa_0,N}\).
Thus mutual information of the same input \(\sigma\) for every original
signal is at least its value for that worst pair.
The cuts, \(U_E,W_E\), and the signal-norm bound in (14),(15) are
uniform with \(\kappa\ge\kappa_0,n\le E\), replacing \(\kappa\) by
\(\kappa_0\) in the norm bound. Apply the finite-family-independent
coding step to the entire signal/reset index set.
For the converse choose the worst pair and its witness.
Do not degrade the full reset mixture: receiver processing could change
the reset entropy and leave the budget class.
\end{proof}

\section{A strict flag-loss gap at bounded signal information}
The following combines established binary-ensemble inequalities.
In particular the purification/Gram bound in its proof is already
explicit in \cite[Corollary~4, Lemma~5]{roga} and
\cite[Eqs.~(3)--(6)]{fannes}; it is not a new lemma claimed here.
Use root fidelity \(F(r,s)=\|\sqrt r\sqrt s\|_1\), rather than its
square, and Bures distance \(d_B(r,s)=\sqrt{2(1-F(r,s))}\).

\begin{lemma}[Strict loss bound]\label{lem:strict}
Let \(\omega_{RB}\) be normalized on separable spaces, with
\(c=I(R:B)_\omega<\infty\). For an arbitrary state \(\tau_B\), put
\(\sigma=\omega_R\otimes\tau\), \(z=\lambda\omega+q\sigma\) and
\(\epsilon_\lambda(x)=2\lambda q\,2^{-x}/(3\ln2)\). Then
\[
 0\le \lambda c-I(R:B)_z
 \le h(\lambda)-\epsilon_\lambda(c).                 \tag{16}
\]
No energy or entropy condition on \(\tau\) is required.
\end{lemma}
\begin{proof}
Adjoin the binary branch flag \(X\). The two branches have the same
reference marginal, so \(I(R:X)=0\), \(I(R:BX)=\lambda c\).
Writing \(\chi_{RB}=I(X:RB)\), \(\chi_B=I(X:B)\), the loss \(L\) obeys
\[
 L=\chi_{RB}-\chi_B,\qquad
 h(\lambda)-L=\chi_B+H(X|RB).                         \tag{17}
\]
All flag quantities are finite. These are chain-rule identities, not
subtractions of possibly infinite reset or output entropies.

The established bound \(D(r\Vert s)\ge-2\log_2 F(r,s)\)
\cite{renyi} gives
\(D(r\Vert s)\ge d_B(r,s)^2/\ln2\).
With \(m=\lambda\omega_B+q\tau\) and \(d=d_B(\omega_B,\tau)\),
weighted Cauchy--Schwarz and the Bures triangle inequality yield
\[
 \chi_B\ge\frac{\lambda d_B(\omega_B,m)^2+
                     q d_B(\tau,m)^2}{\ln2}
          \ge\frac{\lambda q d^2}{\ln2}.              \tag{18}
\]
Let \(f_0=F(\omega,\sigma)\). Uhlmann purifications of these states
have overlap arbitrarily close to \(f_0\). Holevo information cannot
increase under partial trace, whence the known Gram bound is
\[
 \chi_{RB}\le S(G),\qquad
 G=\begin{pmatrix}\lambda&\sqrt{\lambda q}f_0\\
                 \sqrt{\lambda q}f_0&q\end{pmatrix}.
\]
Since \(G\) has diagonal \((\lambda,q)\), qubit Pinsker gives
\[
 H(X|RB)\ge h(\lambda)-S(G)
       =D(G\Vert\operatorname{diag}(\lambda,q))
       \ge\frac{2\lambda q f_0^2}{\ln2}.              \tag{19}
\]
Set \(a=F(\omega,\omega_R\otimes\omega_B)\).
Choose purifications of this pair with overlap approaching \(a\),
and a purification of \(\sigma\) at vector distance approaching \(d\)
from the second one. Their overlaps show
\(f_0\ge(a-d)_+\). Also \(a^2\ge2^{-c}\).
Consequently (17)--(19) and the scalar inequality
\[
 d^2+2(a-d)_+^2\ge\tfrac23a^2
\]
prove (16). For \(d\le a\) the remainder is
\(3(d-2a/3)^2\); for \(d\ge a\) it follows from \(d^2\ge a^2\).

The infinite-dimensional steps do not require finite separate entropies.
For completeness, finite-dimensional \(D\ge-2\log_2 F\) follows by
spectral Jensen, followed by
\(\operatorname{Tr}\sqrt r\sqrt s\le\|\sqrt r\sqrt s\|_1\).
Common finite-rank compression with an overflow flag, data processing
for \(D\), and trace-norm continuity of fidelity give the separable
version, including singular supports by limits. The Gram matrix
always has dimension two, and approximate purifications suffice.
The flag chain rules are well-defined because \(c\) is finite and
the flag is finite. Thus no infinity-minus-infinity expression occurs.
\end{proof}

\begin{proposition}[Operational gap for all fixed resets]\label{prop:gap}
For each fixed \(E>0\) and finite \(n\ge0\), the full compound reset
family, and hence every entropy-budget subfamily, has capacity at least
\[
 \lambda C_{\kappa,n}(E)-h(\lambda)
          +\epsilon_\lambda(C_{\kappa,n}(E)).        \tag{20}
\]
The same bound holds for every known reset.
For unknown fixed signal parameters
\(\kappa\in[\kappa_0,1]\), \(n\in[0,N]\), it holds with
\(C_{\kappa_0,N}(E)\), for every finite \(E,N\), even \(N\gg E\).
\end{proposition}
\begin{proof}
Write \(C_*=C_{\kappa,n}(E)\), or \(C_{\kappa_0,N}(E)\) for the
envelope, and \(G_\lambda(x)=\lambda x+\epsilon_\lambda(x)\).
It is strictly increasing for \(x\ge0\), since
\(G_\lambda'(x)=\lambda-(2\lambda q/3)2^{-x}>0\).
Pointwise application of the state lemma alone is insufficient to
establish a common operational code, so we give that step explicitly.

Fix an error in rate \(\eta>0\). Start with
\(\gamma_{E-s}\), where \(s>0\) leaves physical energy slack.
Condition it on input levels \(0,\ldots,J-1\), giving \(\sigma_J\)
of mean at most \(E-s\). Its purifications converge to those of the
thermal state. Lower semicontinuity of mutual information and thermal
optimality at energy \(E-s\) imply
\(I(\sigma_J,\Phi_*)\to C_*(E-s)\).
Choose \(s,J\) with \(I(\sigma_J,\Phi_*)\ge C_*-\eta/2\).
The signal degradation used in Proposition~\ref{prop:unknown}
compares this same input with every signal in the envelope, for
arbitrary finite \(N\), not just \(N\le E\).

Now apply the physical output cutoff \({\cal D}_T\), retaining levels
\(0,\ldots,T\) and sending the rest to an orthogonal flag.
All signal outputs have mean at most \(E+N\).
Their cutoff tail is at most \(p_T=(E+N)/(T+1)\).
After insertion into a common space the trace distance is at most
\(e_T=\sqrt{p_T}+p_T/2\).
With the fixed reference dimension \(J\), Winter's conditional-entropy
bound \cite[Lemma~2]{winter} bounds the information error by
\[
 2e_T\log_2 J+(1+e_T)h(e_T/(1+e_T))\longrightarrow0.
\]
Thus one finite \(T\) gives signal information at least \(C_*-\eta\)
for all signals. The output entropy of the reset was never bounded.
Each effective mixture is finite-dimensional and has arbitrary reset
\({\cal D}_T(\tau)\). Lemma~\ref{lem:strict} and monotonicity of
\(G_\lambda\) give the same information lower bound
\(G_\lambda(C_*-\eta)-h(\lambda)\) for every channel, taking
\(\eta<C_*\).

For clarity, the cost-preserving type step can be made exact.
Approximate the Fock probabilities of \(\sigma_J\) by a type \(p_l\)
with mean at most \(E\), and let \(\pi_l\) be uniform on that exact
type subspace. Put \(\rho_l=\operatorname{diag}(p_l)\) and
\(d_l=D(\pi_l\Vert\rho_l^{\otimes l})\le J\log_2(l+1)\).
For any effective finite channel \(M\), its output and complementary
total correlations \(D_{\rm out},D_{\rm env}\) satisfy
\[
 I(\pi_l,M^{\otimes l})
 =lI(\rho_l,M)-d_l-D_{\rm out}+D_{\rm env}
 \ge lI(\rho_l,M)-2d_l .
\]
Here \(D_{\rm out}\le d_l\) by data processing; both relative entropies
are finite by domination of the type state by its product state.
Finite reference dimension makes
\(I(\rho_l,M)\to I(\sigma_J,M)\) uniform in \(M\).
Apply the arbitrary-index finite compound coding theorem
\cite[Theorem~1, Lemma~9]{compound} on this same type subspace.
The classical assisted rate is twice its quantum assisted rate by
superdense coding. Every encoded state has mean physical cost at most
\(E\) per use. First take the outer coding limit at fixed cuts and type;
then optimize the finite type, cuts and slack. This proves (20).
The cutoff is implemented at the receiver of the original channel.
No uniform blocklength in \(E,N\), or changing channel within a block,
is asserted.
\end{proof}

\section{Necessity of the sublinear thermal regime}
\begin{proof}[Proof of necessity in Theorem~\ref{thm:sharp}]
For fixed \(r>0\), put \(n=rE\). In (C), \(M/E\to\kappa+\delta r\) and
\[
 D/E\to d_r=\sqrt{[\delta+(1+\kappa)r]^2-4\kappa r^2},
 \qquad
 u_{l,s}/E\to\tfrac12[d_r\mathbin{\pm}\delta(1-r)].
\]
The identity \(d_r^2-\delta^2(1-r)^2=4\delta r>0\)
shows that both complementary occupations grow linearly with
positive coefficients and \(u_lu_s/E^2\to\delta r\).
The expansion \(g(x)=\log_2x+\log_2e+o(1)\) therefore yields
\[
 C_{\kappa,rE}(E)\longrightarrow
 K_r=\log_2\left(1+\frac{\kappa}{\delta r}\right).     \tag{21}
\]
If \(n(E)\) is not \(o(E)\), there are \(r>0\) and an unbounded
subsequence on which \(n(E)\ge rE\).
Receiver Gaussian-noise order and (21) give
\(C_{\kappa,n(E)}(E)\le K_r+1\) eventually there.
Proposition~\ref{prop:gap} implies, uniformly over every budget,
\[
 D_Y(E,n(E),B)
 \le h(\lambda)-\epsilon_\lambda(K_r+1),
 \qquad Y\in\{{\rm ind},\cmp\}.                     \tag{22}
\]
Take \(B(E)=\log_2(E+1)\). The target curve is exactly \(h(\lambda)\)
for Shannon, and tends to it for every fixed finite \(\alpha>1\).
The positive constant in (22) prevents the complete uniform law.
This also covers bounded signal benchmarks larger than the flag
entropy, where nonnegativity of capacity alone would not suffice.

For the envelope use the same reasoning at the worst signal
\((\kappa_0,N(E))\), with the compound version of (20).
Its construction was valid at every finite \(E,N\).
If \(N/E\) fails to vanish there is the same positive subsequential
separation. The preceding sublinear proof, including pure-loss
subsequences, supplies the other implication.
\end{proof}

\section{The complete unassisted thermal curves}
\label{sec:unassisted}
This section changes the operational resource: no preshared entanglement
is supplied. Unrestricted block inputs may be entangled across uses;
the constructive lower bound uses local mixed preparations and a
collective quantum receiver. All other cost, error and fixed-channel
conventions remain as above. Define \(C_{\rm cl}\) to be this unassisted
classical capacity and put
\[
 G_{\kappa,n}(E)=g(\kappa E+t)-g(t),\qquad t=(1-\kappa)n.
\]
This established signal capacity is attained by a Gaussian coherent
ensemble \cite{gaussiancl}. Its multimode minimum-output-entropy theorem
is imported, not reproved here. Replace \(\CEA\) by \(C_{\rm cl}\) and
the benchmark \(C\) by \(G\) in the individual and compound deficits;
write the results \(D^U_Y\). For an unknown signal envelope use
\(G_\star=G_{\kappa_0,N}(E)\).

\begin{theorem}[Complete unassisted thermal curves]
\label{thm:unassisted}
Fix the interior parameters and entropy orders above.
For every nonnegative thermal path \(n(E)\), both
\[
 \sup_{B\ge0}|D^U_Y(E,n(E),B)-f(t_\alpha(E,B))|\longrightarrow0,
 \qquad Y\in\{{\rm ind},\cmp\},
\]
hold if and only if \(n(E)=o(E)\).
For the unknown fixed family
\(\kappa\in[\kappa_0,1]\), \(n\in[0,N(E)]\), the analogous complete
curves hold if and only if \(N(E)=o(E)\).
Zero, bounded and oscillating baths are included. Known \(\kappa\)
and worst-family \(\kappa_0\) stay strictly below one.
There is no uniform blocklength claim.
\end{theorem}

\subsection{Same-input transfer and the Gaussian gap}
Let \(\omega_{RB}\) be the signal state with a fixed finite-entropy
reference \(\rho_R\), and let a measurement \(R\to X\) prepare an
ensemble. Write \(I_s=I(R:B)_\omega\), \(J_s=I(X:B)_\omega\), and use
\(I_z,J_z\) for the corresponding quantities after the reset mixture.
For the branch flag \(F\), independent of \(R,X\), chain rules give
\[
 I_z=\lambda I_s-L_R,\quad J_z=\lambda J_s-L_X,\qquad
 L_R=I(R:F|B),\quad L_X=I(X:F|B).
\]
Data processing is on \(R\), not on the conditioning system \(B\);
thus \(0\le L_X\le L_R\) and
\begin{equation}
 0\le I_z-J_z\le\lambda(I_s-J_s).          \tag{U1}\label{eq:transferU}
\end{equation}
These are finite relative-entropy informations bounded by
\(2S(\rho_R)\); the losses are at most \(h(\lambda)\).
There is no subtraction of infinite reset entropies.
Infinite-dimensional CMI and its data processing are established
\cite[Section6.1, Theorem2]{cmiinf}. For continuous heterodyne outcomes,
use nested finite Borel partitions generating the outcome sigma-algebra.
Relative-entropy information increases to its continuous-ensemble value,
so the inequalities pass to the limit; differential entropy is not used.

For a purification of \(\gamma_E\), reference heterodyne has amplitude
variance \(E+1\) and prepares coherent input amplitude
\(\sqrt{E/(E+1)}\,\overline{\alpha}\), of variance \(E\).
Thus \(I_s=C\), \(J_s=G\). The emitter prepares the resulting coherent
states locally; it need not share or measure the analytical reference.
Equation \eqref{eq:transferU} compares this input and ensemble, not the
optimized assisted and unassisted capacities for arbitrary resets.

For \(E,t>0\), set \(A=\delta(E+n)+1\). Direct algebra gives
\[
 D^2-A^2=4\kappa Et,\qquad
 u_lu_s=Et\left(1-\frac{2\kappa}{D+A}\right),\qquad
 \frac{2\kappa}{D+A}\le\frac1{t+1}.
\]
The expansion
\[
 g(x)=\log_2x+\log_2e+r(x),\quad
 r(x)=\frac1{\ln2}\int_0^1\ln(1+s/x)\,ds,\quad
 0\le r(x)\le\frac1{2x\ln2}
\]
then yields the uniform bound
\begin{equation}
 0\le\Delta_{\kappa,n}(E):=C_{\kappa,n}(E)-G_{\kappa,n}(E)
 \le\log_2(1+1/t)+\frac{1/E+1/t}{2\ln2}.   \tag{U2}\label{eq:gapU}
\end{equation}
Negative remainders from \(g(u_l),g(u_s)\) only improve this upper
bound. Nonnegativity is reference-measurement data processing.
The bound needs \(t>0\) but no \(n=o(E)\); for \(t\ge cE\) it is
\(O(E^{-1})\) uniformly even when \(\kappa\) approaches one.
At \(\kappa=1\), treat \(t=0\) separately:
\(G=g(E)\), \(C=2g(E)\).

\subsection{A common unassisted code with physical cost}
Here are the operational steps behind the lower bound; pointwise
Holevo information alone would not suffice. At fixed finite \(E,N\),
start with the coherent Gaussian ensemble of mean \(E'=E-s<E\).
Condition its amplitude on \(|\beta|^2\le L\). The discarded probability
is \(\varepsilon=e^{-L/E'}\), its conditional mean cost is \(L+E'\),
and the Holevo loss for every output channel is at most
\[
 h(\varepsilon)+\varepsilon g(L+E').
\]
This follows by retaining the disk indicator and applying data processing
to its tail ensemble at the input. The retained average obeys
\(\rho_L\le\gamma_{E'}/(1-\varepsilon)\).
Partition the disk into finitely many cells of diameter \(d\).
Each letter is the conditional mixed coherent preparation in its cell,
preserving the average and cost exactly. After displacement by a cell
representative its mean energy is at most \(d^2\); forgetting its internal
position loses at most \(g(d^2)\) information through any channel.
The letter support need not have finite Fock dimension.

Cut the receiver output at \(T\) with an orthogonal overflow flag.
The signal average energy is at most \(E+N\), so its overflow probability
is at most \((E+N)/(T+1)\). With \(J\) letters, finite-reference
continuity \cite[Lemma2]{winter} controls its Holevo loss uniformly by
\[
 W(e,J)=2e\log_2J+(1+e)h(e/(1+e)),\qquad
 e\le\sqrt{p_T}+p_T/2.
\]
The reset is not approximated uniformly in trace norm. Its original
top-\(m\) part is mapped to a support of rank at most \(m+1\), including
overflow, capturing at least \(1-\eta_m\) reset mass. The compressed
signal mass there is at most \((m+1)\nu\), where
\[
 \nu=\max\left\{\frac1{(1+\kappa_0E')(1-\varepsilon)},p_T\right\}.
\]
The induced flag-entropy bound is
\[
 H(F|B_{\rm cut})\le f(\eta_m)
       +b(\min\{\lambda,\lambda(m+1)\nu\},q).
\]
The support test estimates this entropy; it is not the common decoder.
It uses the reset's original spectral tail, without asserting preservation
of its entropy budget by compression.

Approximate the \(J\) probabilities by a rational type with strict cost
slack. Every word has the same expected physical cost. For a type of
length \(\ell\), uniform type-word preparation has Holevo information
\[
 \chi_{\rm type}=\ell\chi(p_\ell,W)
       -D(\omega_\ell\Vert W(p_\ell)^{\otimes\ell}),\qquad
 D(\omega_\ell\Vert W(p_\ell)^{\otimes\ell})
       \le d_\ell\le J\log_2(\ell+1).
\]
The first identity uses its identical one-site marginals; the inequality
is input type divergence followed by data processing. There is only one
\(d_\ell\) loss. Apply the arbitrary-index finite cq compound coding
theorem \cite[Theorem3.1 and direct Theorem5.10]{cqcompound} to the finite
superalphabet of these words and the finite output. Its unknown index
includes \(\kappa,n,\tau\); encoder and decoder are common. The outer
code concatenates type words and hence preserves the physical cost of
each message. Fix all cuts, letters and types before taking its coding
limit, then optimize these finite choices.

For explicit budget uniformity take \(E'=E-1\), make the disk, cell and
continuity errors at most \(E^{-2}\), and choose \(p_T\le E^{-3}\).
The benchmark cost loss is at most \(\log_2(E/(E-1))\), obtained by
integrating its derivative with signal noise fixed. For Shannon choose
\(m=\lfloor E/\ln(E+1)\rfloor\), and for fixed \(\alpha>1\) choose
\(m=\lfloor\sqrt E\rfloor\). The spectral-tail bounds and continuity of
\(f\) used above then give
\begin{equation}
 C_{{\rm cl},\cmp}\ge\lambda G_\star-f(t_\alpha(E,B))-o(1),
                                                 \tag{U3}\label{eq:cqU}
\end{equation}
uniformly all finite budgets, whenever \(N=o(E)\).
For a known signal use its \(G\). This lower construction itself does
not need a diverging bath.

\subsection{Diverging-bath comparison and necessity}
An unassisted code is an admissible assisted code. Therefore
\[
 D^U_Y\ge D^{\rm EA}_Y-\lambda(C_\star-G_\star).
\]
For known \(n\to\infty\), or the worst envelope \(N\to\infty\),
\eqref{eq:gapU} makes this gap vanish. The earlier assisted curve and
\eqref{eq:cqU} sandwich both unassisted deficits. Only the worst-pair
gap is needed; it is not assumed small in each cold family member.
This supplies an alternative sufficiency proof for a diverging bath.
The block argument below removes that premise; neither route assumes
additivity of the mixed reset channel.

For necessity the strict state inequality above and
\eqref{eq:transferU} give, for the same coherent ensemble,
\[
 J_z\ge\lambda G-h(\lambda)+\epsilon_\lambda(C).
\]
To obtain a uniform code at any fixed finite \(E,N\), not only in a
sublinear envelope, compress the thermal reference analytically onto
levels \(0,\ldots,J_R-1\) with a failure flag. Its tail mass is
\(\varepsilon_R=(E/(E+1))^{J_R}\); its normalized tail entropy is \(g(E)\).
Recording and discarding the binary sector, then collapsing its tail,
costs at most
\[
 A_{J_R}=2h(\varepsilon_R)+2\varepsilon_Rg(E)
\]
mutual information, uniformly for every signal. With that finite
reference, the physical output cutoff has information error at most
\[
 0\le I_i-I_i^T
 \le A_{J_R}+W(\sqrt{p_T}+p_T/2,J_R+1).
\]
First take \(J_R\) large, then \(T\) large, at fixed \(E,N\).
For the coherent ensemble its Holevo loss is at most this information
loss by CMI data processing on \(R\to X\) in an isometric dilation of
the cutoff. Apply the transfer and strict inequalities to the same
effective signal and compressed reset. Subsequent disk, cell and type
steps yield a common operational strict rate. One can start at mean
\(E\): conditioning on a finite disk creates the strict slack for types.
This proves a known-signal deficit at most
\(h(\lambda)-\epsilon_\lambda(C)\) for every budget.
If \(n/E\) fails to vanish, the earlier critical benchmark limit bounds
\(C\) on a subsequence and the gap stays positive.

For an unknown envelope with \(N/E\) nonvanishing along a subsequence,
write \(G_\star\le C_\star\le K<\infty\) there. The common thermal input
has \(G_i\ge G_\star\). Members with \(G_i\ge G_\star+1\) already give
gain at least \(\lambda\) above \(\lambda G_\star-h(\lambda)\).
In the other stratum \(G_i<K+1\), use
\[
 G_i\ge\log_2\left(1+\frac{\kappa_iE}{t_i+1}\right),\qquad
 t_i\ge\frac{\kappa_0E}{2^{K+1}-1}-1 .
\]
Equation \eqref{eq:gapU} gives \(C_i-G_i=O(E^{-1})\) uniformly,
so \(C_i\le K+2\) eventually and the gain is at least
\(\epsilon_\lambda(K+2)>0\). The identity signal lies in the first
stratum. The uniform cutoff argument preserves one common code, not
one code per stratum or member. At the saturating budget
\(B(E)=\log_2(E+1)\), this positive separation excludes both complete
curves as in the assisted proof, proving necessity.

\subsection{A converse uniform in block length}
\label{sec:blockU}
Fix the same finite-energy witness
\(\tau_a=(1-a)|0\rangle\langle0|+a\gamma_M\), \(a\in[0,1]\),
throughout each block. Write
\({\cal N}={\cal D}_{\tau_a}\circ\Phi\), with
\({\cal D}_{\tau_a}=\lambda\id+q\tau_a\Tr\).
Consider any finite ensemble \(\{p_x,\rho_x^{(m)}\}\) of \(m\)-mode
inputs, possibly entangled across uses, of total mean cost at most \(mE\).
Its Holevo information is \(J_{\cal N}^{(m)}\).
All entropies below are finite at each finite \(m,E,n\): labels, input
mean energy and witness energy are finite. Other admissible resets
are not assumed to have finite energy or entropy.

At use \(j\), let \(R_j=XB_{<j}\), with past outputs processed by
\({\cal N}\), and let \(\omega_j\) be this reference and the current
signal output before its reset mixture. Its branch flag \(F_j\) is
independent of \(R_j\). Put \(K_j=S(B_j|R_j)_{\omega_j}\) and
\(H_j=H(F_j|R_jB_j)\). The exact causal entropy chain rule is
\begin{equation}
 J_{\cal N}^{(m)}=S(B^m)_{\rm av}
       -\lambda\sum_jK_j-mqS(\tau_a)-mh(\lambda)+\sum_jH_j .
 \tag{U4}\label{eq:causalU}
\end{equation}
The reference includes the message and quantum past; it is not a
receiver supplied with the message. Degrading the all-signal past by
\({\cal D}_{\tau_a}\) increases conditional entropy. The established
tensor-power minimum-output theorem
\cite[Eq.(2)]{gaussiancl}\cite[Proposition4, Sections6--7]{optimizer} gives
\begin{equation}
 \sum_jK_j\ge S(B^m|X)_{\Phi^{\otimes m}}\ge m g(t).
 \tag{U5}\label{eq:summedU}
\end{equation}
This is a summed floor; individual \(K_j\) may be negative.
The flat Choi order extends from purifications to mixed quantum
references by reference processing. Hence
\[
 \omega_j\le c\rho_{R_j}\otimes I,\quad
 c=(t+1-\kappa)^{-1},\qquad d_j=\log_2c+K_j\ge0.
\]

Blocks with \(J_{\cal N}^{(m)}/m<\lambda G-h(\lambda)-1\) already
satisfy the desired converse. Consider all others.
Let \(\mu_j,S_j\) be the mean and entropy of the average signal marginal.
The cost gives \(\sum_j\mu_j\le mM\), not an individual bound
\(\mu_j\le M\). Group signal and thermal reset against vacuum, of weight
\(P_a=\lambda+qa\), then use concavity after summing positions:
\[
 S(B^m)_{\rm av}\le m[h(P_a)+P_ag(M)],\quad
 S(\tau_a)\ge h(a)+ag(M)-h(a/(M+1)).
\]
Equation \eqref{eq:causalU} therefore yields
\begin{equation}
 J_{\cal N}^{(m)}
 \le m\lambda G-mf(a)-\lambda\left(\sum_jK_j-mg(t)\right)
              +\sum_jH_j+mqh(a/(M+1)).          \tag{U6}\label{eq:blockboundU}
\end{equation}
Using \(H_j\le h(\lambda)\) and the near-capacity assumption gives
\[
 \frac1m\sum_jd_j\le C_d,\qquad
 C_d=\log_2e-\log_2(1-\kappa)
                  +\frac{2h(\lambda)+1+q}{\lambda}.
\]
Here \(g(t)+\log_2c\le\log_2e-\log_2(1-\kappa)\).
The constant is independent of block length, noise, reset weight,
message count and energy sharing.

The alternative bound
\(S(B^m)_{\rm av}\le\sum_j[\lambda S_j+qS(\tau_a)+h(\lambda)]\)
shows that \(\sum_jS_j\ge mg(M)-mC_s\),
\(C_s=(2h(\lambda)+1)/\lambda\).
With the common thermal comparator \(\gamma_M\) and
\(\beta_M=\log_2(1+1/M)\),
\[
 \sum_jD((\omega_j)_{B_j}\Vert\gamma_M)
 =mg(M)-\sum_jS_j+\beta_M\left(\sum_j\mu_j-mM\right)\le mC_s.
\]
Binary vacuum/rest data processing bounds its average vacuum mass by
\[
 \frac1m\sum_j w_j\le\frac{C_s+1}{\log_2(M+1)}.
\]
No individual near-thermal output assumption is needed.

Set \(S=c(M+1)\), \(u=\sqrt{\ln S}\),
\(T=\lceil(M+1)u\rceil\), \(P_B=\sum_{r=1}^T|r\rangle\langle r|\),
\(p_j=\Tr(I_{R_j}\otimes P_B)\omega_j\), \(v=\Tr P_B\gamma_M\).
For \(n=o(E)\), \(M,S\to\infty\) and \(v\to1\). The total-energy
tail bound gives
\[
 \frac1m\sum_j(1-p_j)\le
 A_E:=\frac{C_s+1}{\log_2(M+1)}+\frac1u\longrightarrow0.
\]
Let \(\ell_E=\log_2S\), \(\varepsilon_E=\sqrt{A_E}\).
Call a position good when \(p_j\ge1-\varepsilon_E\) and
\(d_j\le\sqrt{\ell_E}\). Markov bounds the fraction of bad positions:
\begin{equation}
 \frac{\#\{\text{bad }j\}}m
 \le\sqrt{A_E}+\frac{C_d}{\sqrt{\ell_E}}\longrightarrow0.
 \tag{U7}\label{eq:badU}
\end{equation}
This uses no union bound over the block.

For each good position use the original reference \(\rho_{R_j}\)
in the earlier pinching argument. With the normalized window state
\(\omega_{jg}\) and
\(\sigma_{jg}=\rho_{R_j}\otimes P_B\gamma_MP_B/v\), it gives
\[
 D_0(\omega_{jg})\ge\log_2c-\frac{d_j+h(p_j)}{p_j},\quad
 L_j=(S/p_j)v(M/(M+1))^{-T},\quad \omega_{jg}\le L_j\sigma_{jg},
\]
\[
 D(\omega_{jg}\Vert\sigma_{jg})
 \ge\ell_E+\log_2v-\frac{d_j+1}{p_j},\quad
 \log_2L_j=\ell_E+\log_2v-\log_2p_j+T\beta_M .
\]
The functional \(D_0\) includes
\(D(\rho_{R_j}^g\Vert\rho_{R_j})\); it is not posterior conditional
entropy alone. Common finite compressions with a normalized overflow
flag preserve order and converge in trace norm. Data processing and
lower semicontinuity extend normalized TRE to this separable reference
with the same dimension-independent constant.

An explicit error is obtained by defining
\[
 D_-=\ell_E+\log_2v-\frac{\sqrt{\ell_E}+1}{1-\varepsilon_E},
 \quad L_+=\ell_E-\log_2(1-\varepsilon_E)+T\beta_M,\quad
 r_E=1-\frac{(1-1/(Sv))D_-}{L_+}.
\]
These formulas are used eventually when \(D_->0\), \(Sv>1\)
and \(\varepsilon_E<1/2\); \(r_E\to0\).
TRE gives \(\tfrac12\|\omega_{jg}-\sigma_{jg}\|_1\ge1-r_E\).
A Helstrom test supported on the nonvacuum window has summed normalized
errors at most \(r_E\). On the original signal its failure is at most
\(\varepsilon_E+r_E\); on the reset its acceptance is
\(av\Tr Q_j\sigma_{jg}\). The binary branch error is at most
\(\varepsilon_E+r_E\). For large energy Fano and \eqref{eq:badU} yield
\begin{equation}
 \frac1m\sum_jH_j
 \le h(\varepsilon_E+r_E)
       +h(\lambda)\left(\sqrt{A_E}+\frac{C_d}{\sqrt{\ell_E}}\right)
 =:\eta_E\longrightarrow0.                    \tag{U8}\label{eq:flagmeanU}
\end{equation}
The tests bound already-conditioned entropies, not an operational code.

Equations \eqref{eq:summedU}, \eqref{eq:blockboundU} and
\eqref{eq:flagmeanU} now imply for every block ensemble and length
\[
 J_{\cal N}^{(m)}/m\le
       \lambda G-f(a)+qh(a/(M+1))+\eta_E.
\]
The error is uniform in length, message count, entanglement, energy sharing
and \(a\). Holevo--Fano bounds the regularized operational capacity by
the same expression, taking its coding limit before the energy limit.
The admissible Shannon witness is the clipped choice from Section6;
for fixed \(\alpha>1\), choose \(a=t_\alpha\) exactly, since
\(\Tr\tau_a^\alpha\ge(1-a)^\alpha\).
Together with \eqref{eq:cqU}, this proves sufficiency without bath
divergence. For the unknown family the converse selects only its worst
signal; the lower code remains common. The preceding necessity finishes
both equivalences, without assuming additivity of the mixed channel.

\section{Amplification, additive noise and a larger unknown family}
\label{sec:generalU}
The unassisted argument has a wider signal class. Let \(\Psi_{\kappa,t}\)
be a centered phase-insensitive gauge-covariant Gaussian channel with
\[
 \chi_{\rm out}(z)=\chi_{\rm in}(\sqrt\kappa z)
       e^{-[t+(1-\kappa)/2]|z|^2},\qquad
 \kappa>0,\quad t\ge(\kappa-1)_+ .
\]
Its vacuum output is \(\gamma_t\), its output mean for input mean \(e\)
is \(\kappa e+t\), and its established unassisted signal capacity is
\(G_{\kappa,t}(E)=g(\kappa E+t)-g(t)\)
\cite[Eqs.(24)--(29)]{gaussiancl}\cite[Section7]{optimizer}.
The tensor-power minimum output entropy is \(m g(t)\), including
entangled finite-energy inputs.
For \(0<\kappa<1\), \(t=(1-\kappa)n\) is thermal loss;
for \(\kappa>1\), \(t=(\kappa-1)(n+1)\) is thermal amplification;
for \(\kappa=1\), \(t\) is additive classical Gaussian noise.
These Gaussian formulas are antecedents, not new signal-capacity claims.
Use the same reset mixture and unassisted operational conventions,
and define its deficits relative to \(\lambda G_{\kappa,t}\).
For fixed \(0<\kappa_0<1\le\kappa_+<\infty\), write
\[
 {\cal F}_{E,T,B}=
 \left\{\lambda\Psi_{\kappa,t}+q\tau\Tr:
  \kappa_0\le\kappa\le\kappa_+,\ 
  (\kappa-1)_+\le t\le T,\ \tau\in{\cal T}_{\alpha,B}\right\}.
\]
The two unknown-family deficits mean, explicitly,
\[
\begin{split}
 D^U_{{\rm ind},{\cal F}}(E,T,B)
 &=\lambda G_{\kappa_0,T}(E)
       -\inf_{{\cal M}\in{\cal F}_{E,T,B}}C_{\rm cl}({\cal M},H,E),\\
 D^U_{\cmp,{\cal F}}(E,T,B)
 &=\lambda G_{\kappa_0,T}(E)
       -C_{{\rm cl},\cmp}({\cal F}_{E,T,B},H,E).
\end{split}
\]
The infimum in the first line allows a different optimal code for every
fully known channel, including its reset. The second line requires one
encoder and decoder, with uniform error over the joint unknown index.
All coding limits are taken at fixed finite \(E,T,B\).

\begin{theorem}[Complete curves for covariant Gaussian signals]
\label{thm:generalU}
For every fixed \(\kappa\ne1\), the two complete entropy-budget
curves equal \(f(t_\alpha)+o(1)\), uniformly all finite budgets,
if and only if \(t(E)=o(E)\).
For fixed \(\kappa=1\), the same equivalence holds under an eventual
lower bound \(t(E)\ge t_0>0\).

Fix \(0<\kappa_0<1\) and finite \(\kappa_+\ge1\). For the unknown
fixed family of all physically admissible pairs
\[
 \kappa\in[\kappa_0,\kappa_+],\qquad
 (\kappa-1)_+\le t\le T(E),
\]
the corresponding two deficit curves relative to
\(\lambda G_{\kappa_0,T}(E)\) hold if and only if \(T(E)=o(E)\).
This includes bounded, zero and oscillating envelopes. Pure loss and
identity may be members; amplifiers are included whenever the physical
noise bound permits them.
\end{theorem}

\begin{proof}
We give the exact changes to the preceding proof rather than assume
that all Gaussian channels have the same operational constraints.
The established Choi/reference order is
\[
 \omega_{RB}\le c\,\rho_R\otimes I,\qquad
 c=\frac1{t+(1-\kappa)_+}.
\]
For \(\kappa\ne1\), Holevo's amplitude is \(\sqrt\kappa\) and noise
coefficient is \(m_H=t+(1-\kappa)/2\); his norm is
\([m_H+|1-\kappa|/2]^{-1}\)
\cite[Proposition2 and p.13 example]{choi}.
For \(\kappa=1,t>0\), use that source's separate Case3 on p.11:
\(\Delta_K=0,\mu=tI\), with norm \(1/t\).
Its nondegenerate-symplectic proposition is not invoked at \(\kappa=1\).

The averaged-entropy argument of Section~\ref{sec:blockU} only requires
this order, the linear mean, the thermal average and the tensor-power
minimum-output theorem. Its constant \(C_d\) is replaced by
\[
 \frac{2h(\lambda)+1+q}{\lambda}+
 \begin{cases}
  \log_2e-\log_2(1-\kappa),&\kappa<1,\\
  \log_2e+\log_2(1+1/(\kappa-1)),&\kappa>1,\\
  \log_2e+\log_2(1+1/t_0),&\kappa=1.
 \end{cases}
\]
These upper-bound \(g(t)+\log_2c\) using
\(g(t)\le\log_2[e(t+1)]\) and the physical noise floors.
They are fixed along each allowed path. When \(t=o(E)\),
\(M=\kappa E+t\to\infty\) and \(S=c(M+1)\to\infty\).
Thus the same good-position selection and uniform-in-length converse
apply to the fixed witness. The Shannon normalization still holds since
\(g(M)/\log_2(E+1)\to1\).

For the lower code, a truncated coherent Gaussian average satisfies
\(\rho_L\le\gamma_{E'}/(1-\varepsilon)\), so every member's average
signal norm is at most \(1/[(1+\kappa_0E')(1-\varepsilon)]\).
The output mean is at most \(\kappa_+E+T\).
These replace the two signal-specific constants in the disk, cell,
output-cut and exact-type cq construction. The same coherent ensemble
has signal Holevo \(G_{\kappa,t}\ge G_{\kappa_0,T}\):
the formula increases with \(\kappa\) and decreases with \(t\).
This is an information comparison of signals, not an order imposed on
their reset mixtures. For the unknown-family converse choose only the
allowed worst-pair attenuator \((\kappa_0,T)\).

For necessity a different standard Gaussian decomposition avoids a new
assisted-capacity radical formula:
\[
 \Psi_{\kappa,t}={\cal A}_{t+1,0}\circ{\cal L}_{\eta,0},
 \qquad \eta=\frac{\kappa}{t+1}\le1 .
\]
The quantum-limited amplifier has gain \(t+1\) and vacuum noise \(t\);
gain multiplication and characteristic noise recover exactly
\(\Psi_{\kappa,t}\). For the thermal input,
signal-output processing bounds its reference information by the
pure-loss assisted expression
\[
 I_s\le g(E)+g(\eta E)-g((1-\eta)E)
       \le g(\eta E)-\log_2(1-\eta),\qquad \eta<1 .
\]
If \(t/E\ge r>0\) on a subsequence, \(\eta E\le\kappa/r\) and
\(\eta=O(E^{-1})\), so \(I_s\le g(\kappa/r)+1\) eventually.
The strict state gain and the common cq transfer then separate the
deficit from \(h(\lambda)\) at the saturating budget.

For the unknown family on \(T/E\ge r\), \(G_\star\) is bounded by
\(\log_2(1+\kappa_0/r)\). Stratify as before.
The high-\(G_i\) stratum gives gain at least \(\lambda\).
In the other stratum \(G_i\le K+1\), use
\[
 G_i\ge\log_2(1+\kappa_iE/(t_i+1))
       =\log_2(1+\eta_iE).
\]
Thus \(\eta_iE\le2^{K+1}-1\), \(\eta_i=O(E^{-1})\), and
\(I_i\le g(2^{K+1}-1)+1\) uniformly eventually.
The uniform finite-noise cutoff proof applies at each fixed \(E,T\),
with mean bound \(\kappa_+E+T\); it supplies one common strict code
over both strata. This proves necessity without a bath-divergence or
EB premise. It does not determine the critical loss curve.
\end{proof}

The additive boundary \(t\to0\), including an exactly noiseless
\(\kappa=1,t=0\) signal by itself, is not covered by this particular
uniform-order argument. Nor does the theorem include phase-sensitive
squeezing, unbounded gain, unknown phase, or within-block channel
variation. No additivity or Gaussian optimizer is asserted for the
generally non-Gaussian reset mixture.

\section{Evidence, dependencies and remaining problems}
Several nearby model classes precede this work. Finite-dimensional
identity/replacer combinations appear in
Giovannetti--Lloyd--Ruskai \cite[SectionIV.B, Eq.(34)]{glr}.
Convex combinations of bosonic pure-loss channels, including the
loss/vacuum endpoint, are studied by Catalano et al.\ \cite{fading}.
Energy-constrained Gaussian cq compound channels with unknown noise,
phase or attenuation are treated by Cacioppo--N\"otzel--Rosati
\cite{energycompound}. These contextual antecedents are distinct from
claiming new Gaussian signal formulas or a new compound coding theorem.
The editorial abstract of Baitenov \cite{baitenov} concerns a one-shot
capacity of shifted depolarizing channels. Its full text was not obtained;
detailed overlap remains unverified. Literature comparison here is
directed and does not certify worldwide priority.

The strict-gap constant is not claimed optimal. The actual
critical-regime minimax curve, quantitative uniform blocklength bounds
and arbitrary within-block variation remain outside this result.
Neither (21) nor (22) determines the critical loss curve.

The current numerical evidence consists of 90 high-precision scalar
controls for the earlier \(o(\sqrt E)\) working argument and 72
noncommuting finite-state controls for the pinching/order bookkeeping.
The latter test generic finite matrices, not the infinite Gaussian Choi
form or an operational capacity. Neither set proves the asymptotic
theorem. Tests, frozen inputs and reports are retained separately.
An additional 1080 generic finite-state controls of the strict-loss
inequalities were reproduced byte for byte during technical audit627.
They corroborate only finite diagnostics; the state and operational
arguments above are the proof, conditional on the cited dependencies.
The earlier supplement provides three programs: check612 (90 scalar
cases), check614 (72 matrix cases), and check625c (1080 strict-gap
cases), with expected outputs, dependencies and replay records.
The unassisted working argument adds check641 (540 finite-state and
150 Gaussian cases); its supplied outputs were reproduced during
technical audit644. These computations do not prove the cq theorem,
the continuum limit or the Gaussian optimizer.
The new check653 program tests 120 finite two- and three-use causal
chains and a negative individual conditional-entropy example.
The check658 program tests 360 Gaussian parameter bounds and an exact
loss/amplifier characteristic-function identity.
Neither program proves a block converse or capacity. The six programs
and their expected outputs are enumerated in the supplement.
The package does not implement an operational code or certify an
infinite-dimensional optimizer.

\appendix
\section{Model assessments and review context}
The named author directed substantial OpenAI Codex assistance in
derivation, writing and numerical work.
Requested GPT-6 Astra xhigh working-argument audits supported the
fixed-temperature and \(o(\sqrt E)\) versions with their stated limits.
They used shared files and reused context.
Audit615 of the frozen broader working argument614 has since closed
with mathematical support and minor attribution/notation clarifications,
including the Choi form, fixed-reference pinching, infinite-reference
limit, full two-sided minimax bounds and unknown-signal extension.
This revision incorporates those clarifications and the exact Holevo
Choi antecedent. That audit did not assess this exact PDF, certify
importance or independently prove every imported coding theorem.
Full requested-model review620 examined the preceding exact eight-page
draft and supported its main results with a minor correction to the
pinpoint type-bound attribution. This revision corrects that attribution
and supplies the type and input-tail derivations explicitly.
The original report and its objections are retained.
It used reused context and shared files and disclosed incidental exposure
to an older paper's public assessment during a literature search;
no blindness or independence is claimed.
No full assessment or numerical grade of this revised exact PDF exists.
Full review634 of the earlier eleven-page PDF supported its assisted
results, with optional inventory and historical-citation suggestions.
It did not assess the unassisted section added here. This revision
clarifies the program inventory; the earlier binary purification
antecedent of Jain and Nayak \cite[Lemma3.7, Corollary3.8]{jainnayak}
is also credited alongside the existing Gram references.
Technical audit627 of the frozen strict-gap argument supported its
state inequality, common-code transfer at finite \(E,N\), critical
signal limit and necessity for the complete curves.
It assigned no score, did not assess this combined manuscript and did
not certify priority. It reused context from earlier technical audits,
shared files and disclosed earlier assessment exposure; runtime
independence was not attested. Its original report is retained.
Technical audit644 supported the unassisted working argument641,
including the finite-noise uniform cutoff and exact physical-cost
cq code. It imports the earlier assisted and strict-gap arguments;
it is not a full assessment of this exact combined PDF and assigned
no numerical grade. Shared context and prior participation remain
declared. Intermediate audit649 supported a bounded-bath EB witness route;
its original scope and report remain retained. Technical audit654
subsequently supported the thermal all-block argument653 without EB
or bath divergence, including its mixed quantum references and
dimension-independent error uniform in block length.
It assigned no grade and did not assess this combined PDF.
The amplification/additive-noise and enlarged-family section is a
further authorial extension658; it was not audited in654 and must be
examined in the full assessment of this exact manuscript.
The preceding public entropy/optical manuscript has its own preserved
assessment; its grade is not transferred to this draft.
No independently attested runtime isolation or model identity,
independent human review, Lean verification, worldwide originality,
editorial acceptance, or predicted importance is claimed.
Literature reading has been directed: specified Gaussian/coding
passages, Shirokov's mixture bound, and Audenaert's Theorem4, Corollary1
and selected proof pages. It is not an exhaustive literature sweep
or an audit of all dependencies.

\begin{thebibliography}{99}\small
\bibitem{glr} V. Giovannetti, S. Lloyd and M. B. Ruskai,
\emph{Conditions for the multiplicativity of maximal \(l_p\)-norms of
channels for fixed integer \(p\)},
\href{https://arxiv.org/abs/quant-ph/0408103v3}
{arXiv:quant-ph/0408103v3} (2005), SectionIV.B, Eq.(34).
\bibitem{fading} G. Catalano, M. Fanizza, F. A. Mele,
G. De Palma and V. Giovannetti,
\emph{Convex combinations of bosonic pure-loss channels},
\href{https://arxiv.org/abs/2604.26874v2}{arXiv:2604.26874v2}
(2026), selected model and optical endpoints.
\bibitem{energycompound} A. Cacioppo, J. N\"otzel and M. Rosati,
\emph{Compound Channel Capacities under Energy Constraints and Application},
\href{https://arxiv.org/abs/2105.04274v1}{arXiv:2105.04274v1}
(2021), contextual energy-constrained cq coding antecedent.
\bibitem{baitenov} E. L. Baitenov,
\emph{On Jensen Gap and Capacity of a Shifted Depolarizing Quantum Channel},
Proc. Steklov Inst. Math.\ \textbf{324}, 33--43 (2024),
\href{https://doi.org/10.1134/S0081543824010048}
{doi:10.1134/S0081543824010048}. Editorial abstract only;
full text not obtained.
\bibitem{optimizer} V. Giovannetti, A. S. Holevo and R. Garc\'ia-Patr\'on,
\emph{A solution of the Gaussian optimizer conjecture},
\href{https://arxiv.org/abs/1312.2251v1}{arXiv:1312.2251v1}
(2013), Proposition4, Eqs.(25)--(26), Sections6--7; finite second moments.
\bibitem{gaussiancl} V. Giovannetti, R. Garc\'ia-Patr\'on,
N. J. Cerf and A. S. Holevo,
\emph{Ultimate communication capacity of quantum optical channels by solving the Gaussian minimum-entropy conjecture},
\href{https://arxiv.org/abs/1312.6225v1}{arXiv:1312.6225v1}
(2013), selected capacity derivation and Eq.(25).
\bibitem{cqcompound} I. Bjelakovi\'c and H. Boche,
\emph{Classical Capacities of Averaged and Compound Quantum Channels},
\href{https://arxiv.org/abs/0710.3027v2}{arXiv:0710.3027v2}
(2009 revision; first submitted 2007), Theorem3.1 and direct Theorem5.10.
\bibitem{cmiinf} M. E. Shirokov,
\emph{Measures of quantum correlations in infinite-dimensional systems},
\href{https://arxiv.org/abs/1506.06377v2}{arXiv:1506.06377v2},
Section6.1, Theorem2 and selected extension proof.
\bibitem{jainnayak} R. Jain and A. Nayak,
\emph{Accessible versus Holevo Information for a Binary Random Variable}
(2007), Lemma3.7 and Corollary3.8.
\href{https://www.math.uwaterloo.ca/~anayak/papers/JainN07.pdf}
{Author-hosted manuscript}.
\bibitem{hw} A.~S. Holevo and R.~F. Werner,
\emph{Evaluating capacities of Bosonic Gaussian channels},
\href{https://arxiv.org/abs/quant-ph/9912067v1}
{arXiv:quant-ph/9912067v1} (1999), selected Sections IV.C, V.A.
\bibitem{hs} A.~S. Holevo and M.~E. Shirokov,
\emph{On the entanglement-assisted classical capacity of
infinite-dimensional quantum channels},
\href{https://arxiv.org/abs/1210.6926v2}{arXiv:1210.6926v2}
(2013), Theorem1.
\bibitem{gain} A.~S. Holevo,
\emph{The entropy gain of infinite-dimensional quantum channels},
\href{https://arxiv.org/abs/1003.5765v1}{arXiv:1003.5765v1}
(2010), Proposition2.
\bibitem{winter} A. Winter,
\emph{Tight uniform continuity bounds for quantum entropies:
conditional entropy, relative entropy distance and energy constraints},
\href{https://arxiv.org/abs/1507.07775v6}{arXiv:1507.07775v6}
(2016), Lemma2.
\bibitem{compound} M. Berta, H. Gharibyan and M. Walter,
\emph{Entanglement-assisted capacities of compound quantum channels},
\href{https://arxiv.org/abs/1603.02282v2}{arXiv:1603.02282v2}
(2017), Theorem1, Lemma9 and assisted-capacity equivalence.
\bibitem{shirokov} M.~E. Shirokov,
\emph{Convergence conditions for the quantum relative entropy and other
applications of the generalized quantum Dini lemma},
\href{https://arxiv.org/abs/2205.09108v2}{arXiv:2205.09108v2}
(2022), selected Eq.(11), Section5.2.
\bibitem{choi} A.~S. Holevo,
\emph{On the Choi-Jamiolkowski Correspondence in Infinite Dimensions},
\href{https://arxiv.org/abs/1004.0196v3}{arXiv:1004.0196v3}
(2011), Proposition2 and one-mode example on p.13.
\bibitem{audenaert} K.~M.~R. Audenaert,
\emph{Telescopic Relative Entropy},
\href{https://arxiv.org/abs/1102.3040v2}{arXiv:1102.3040v2}
(2011), Theorem4 and Corollary1.
\bibitem{previous} L. Eriksson,
\emph{Sharp entropy budgets for communication with unknown resets and
optical loss}, AIRR.SCIENCE ARR-2026-7SWAJZ0X278GJ9NK, v1
(2026), Working paper.
\url{https://airr.science/papers/ARR-2026-7SWAJZ0X278GJ9NK/versions/v1/}.
\bibitem{roga} W. Roga, M. Fannes and K. \.{Z}yczkowski,
\emph{Universal bounds for the Holevo quantity, coherent information
and the Jensen--Shannon divergence},
\href{https://arxiv.org/abs/1004.4782v1}{arXiv:1004.4782v1}
(2010), Corollary4 and Lemma5.
\bibitem{fannes} M. Fannes, F. de Melo, W. Roga and K. \.{Z}yczkowski,
\emph{Matrices of fidelities for ensembles of quantum states and
the Holevo quantity},
\href{https://arxiv.org/abs/1104.2271v1}{arXiv:1104.2271v1}
(2011), Eqs.(3)--(6); that source uses squared fidelity.
\bibitem{renyi} M. M\"uller-Lennert, F. Dupuis, O. Szehr,
S. Fehr and M. Tomamichel,
\emph{On quantum R\'enyi entropies: a new generalization and some
properties},
\href{https://arxiv.org/abs/1306.3142v4}{arXiv:1306.3142v4}
(2013), order-one-half fidelity identity and Theorem7.
\end{thebibliography}
\end{document}
